% !TEX program = xelatex
\documentclass[12pt,a4paper]{article}

\usepackage[UTF8,scheme=chinese]{ctex}
\usepackage{geometry}
\usepackage{booktabs}
\usepackage{hyperref}
\usepackage{listings}
\usepackage{xcolor}
\usepackage{enumitem}

\geometry{left=2.5cm,right=2.5cm,top=2.8cm,bottom=2.8cm}

\hypersetup{colorlinks=true, linkcolor=blue, urlcolor=blue}

\lstset{
    basicstyle=\ttfamily\small,
    breaklines=true,
    frame=single,
    backgroundcolor=\color{gray!8},
    numbers=none,
    tabsize=4,
    language=Python
}

\title{\textbf{地理围栏实战项目报告} \\ \large 综合实战 · 简易机载逻辑}
\author{实习最终阶段成果 · UAV Communication}
\date{2026年7月13日}

\begin{document}

\maketitle
\tableofcontents
\newpage

\section{项目概述}

\subsection{项目目的}

综合实战项目为\textbf{简易机载逻辑}：实时监测无人机位置，当进入指定圆形地理围栏区域时，自动悬停并反向飞出至围栏边界。

本项目整合前序阶段的 MAVLink 连接、位置解析、模式切换（LOITER/GUIDED）与导航命令（\texttt{DO\_REPOSITION}）。

\subsection{交付物}

\begin{table}[htbp]
\centering
\caption{项目交付物}
\begin{tabular}{p{4.5cm}p{8.5cm}}
\toprule
\textbf{交付物} & \textbf{说明} \\
\midrule
\texttt{geofence\_monitor.py} & 围栏监测与自动响应主脚本 \\
\texttt{geofence\_output\_20260713\_100424.txt} & SITL 验证日志 \\
\texttt{README.md} & 项目说明与运行步骤 \\
本报告 & \texttt{地理围栏实战项目报告.pdf} \\
\bottomrule
\end{tabular}
\end{table}

\section{围栏设计}

\begin{table}[htbp]
\centering
\caption{围栏参数设定}
\begin{tabular}{ll}
\toprule
\textbf{参数} & \textbf{值} \\
\midrule
圆心 & Home 点正北 100\,m \\
半径 & 30\,m \\
进入边界 & Home 北 70\,m 处 \\
退出点 & 圆心正南 30\,m（Home 北 70\,m，南边界） \\
测试北飞目标 & Home 北 115\,m（围栏内部） \\
飞行高度 & 10\,m（相对 Home） \\
\bottomrule
\end{tabular}
\end{table}

从 Home 向北飞行，在距圆心 $\le 30$\,m 时触发响应；反向飞出至南边界（距圆心 30\,m 的南侧点）并悬停。

\section{实现原理}

\subsection{核心流程}

\begin{enumerate}[leftmargin=2em]
    \item 连接 SITL，等待 GPS/EKF 就绪；
    \item 解锁 $\to$ GUIDED $\to$ \texttt{NAV\_TAKEOFF} 至 10\,m；
    \item \texttt{DO\_REPOSITION} 向北飞，循环读取 \texttt{GLOBAL\_POSITION\_INT}；
    \item 进入围栏：LOITER 悬停 3\,s；
    \item GUIDED 飞向退出点，到达后 LOITER 边界悬停 8\,s。
\end{enumerate}

\subsection{围栏判定}

采用平面近似水平距离：

\begin{lstlisting}
def in_geofence(lat, lon, center_lat, center_lon, radius_m):
    return horizontal_distance_m(lat, lon, center_lat, center_lon) <= radius_m
\end{lstlisting}

\subsection{关键 API}

\begin{table}[htbp]
\centering
\caption{主要 MAVLink 接口}
\begin{tabular}{p{4.5cm}p{8.0cm}}
\toprule
\textbf{功能} & \textbf{接口} \\
\midrule
位置监听 & \texttt{recv\_match(type="GLOBAL\_POSITION\_INT")} \\
模式切换 & \texttt{set\_mode\_apm("LOITER")} / \texttt{"GUIDED"} \\
起飞 & \texttt{MAV\_CMD\_NAV\_TAKEOFF} \\
导航 & \texttt{MAV\_CMD\_DO\_REPOSITION}（COMMAND\_INT） \\
GPS 就绪 & \texttt{GPS\_RAW\_INT.fix\_type >= 3} \\
\bottomrule
\end{tabular}
\end{table}

\section{验证结果}

2026-07-13 日志 \texttt{geofence\_output\_20260713\_100424.txt} 摘录：

\begin{lstlisting}
[准备] GPS/EKF 已就绪，可解锁并切换 GUIDED
[命令] TAKEOFF 10.0m ACK | result=ACCEPTED
[监测] 距圆心=39.5m, 围栏内=否
[触发] 进入围栏！| 距圆心=29.3m <= 半径30m
[状态] 当前模式: LOITER
[状态] 到达 围栏南边界 | dist=1.9m
[完成] 围栏触发响应流程结束，飞机在南边界悬停
\end{lstlisting}

全流程耗时约 47 秒，在 SITL 中完成：北飞进入围栏 $\to$ 悬停 $\to$ 反向至南边界 $\to$ 边界悬停。

\section{问题与解决方案}

\begin{table}[htbp]
\centering
\caption{问题排查记录}
\begin{tabular}{p{3.2cm}p{4.8cm}p{5.5cm}}
\toprule
\textbf{现象} & \textbf{原因} & \textbf{处理} \\
\midrule
GUIDED 切换失败 & GPS/EKF 未就绪 & 增加 \texttt{wait\_for\_navigation\_ready()} \\
\texttt{requires position} & SITL 刚启动 & 预热 8\,s，等待 \texttt{fix\_type >= 3} \\
PreArm: EKF3 waiting & 同上 & SITL 启动后多等 15--20\,s \\
未进入围栏 & 北飞目标在围栏外 & 调整 \texttt{--fly-target-north} \\
与 QGC 冲突 & 双端同时发令 & 运行脚本时断开 QGC \\
\bottomrule
\end{tabular}
\end{table}

\section{总结}

本项目实现了以 Home 北 100\,m 为圆心、半径 30\,m 的圆形围栏监测逻辑。无人机从 Home 向北飞行进入围栏后，自动执行悬停、反向飞出至南边界并悬停，在 ArduPilot SITL 中端到端验证通过，达到了综合实战项目的预期目标。

\end{document}
